% Part: computability
% Chapter: computability-theory
% Section: coding-computations

\documentclass[../../../include/open-logic-section]{subfiles}

\begin{document}

\olfileid{cmp}{thy}{cod}
\olsection{کدگذاری محاسبات}

در هر مدل محاسبه می‌توان کارهای زیر را انجام داد:
\begin{enumerate}
\item 
  \emph{تعریف‌های} توابع محاسبه‌پذیر را به شیوه‌ای نظام‌مند وصف کرد.
  برای نمونه، می‌توان مشخصات ماشین‌های تورینگ، تعریف‌های بازگشتی یا
  برنامه‌هایی در یک زبان برنامه‌نویسی را فراهم‌کنندهٔ این تعریف‌ها
  دانست.
\item سابقهٔ کامل محاسبهٔ تابعی را که با تعریفی معین روی ورودی معینی
  داده شده است، وصف کرد. برای نمونه، محاسبهٔ یک ماشین تورینگ را می‌توان
  با دنبالهٔ پیکربندی‌ها (حالت ماشین و محتوای نوار) در هر گام محاسبه
  وصف کرد.
\item آزمود که آیا سابقه‌ای ادعایی از یک محاسبه واقعاً سابقهٔ چگونگی
  محاسبهٔ تابعی محاسبه‌پذیر با تعریفی معین روی ورودی معینی است یا نه
  (ورودی‌ای که تابع روی آن تعریف شده است؛ یعنی محاسبه متوقف می‌شود).
\item از چنین وصفی از سابقهٔ کامل محاسبه، ارزش تابع را برای ورودی معین
  استخراج کرد. برای نمونه، محتوای نوار در آخرین گام یک محاسبهٔ
  متوقف‌شوندهٔ ماشین تورینگ همان ارزش است.
\end{enumerate}

با استفاده از کدگذاری می‌توان به هر وصف از یک تابع محاسبه‌پذیر
\emph{نمایه‌ای} عددی نسبت داد، به‌گونه‌ای که دستورها به روشی
محاسبه‌پذیر از نمایه بازیابی شوند. به همین ترتیب، سابقهٔ کامل یک
محاسبه را نیز می‌توان با عددی واحد کدگذاری کرد. در این صورت رابطهٔ
حسابی «$s$ سابقهٔ محاسبهٔ تابعی با نمایهٔ~$e$ روی ورودی~$x$ را کد
می‌کند» و تابع «خروجی دنبالهٔ محاسبه با کد~$s$» محاسبه‌پذیرند؛ در واقع
هر دو بازگشتی اولیه‌اند.

این واقعیت بنیادی بسیار نیرومند است و اجازه می‌دهد چندین نتیجهٔ مهم و
چشمگیر دربارهٔ محاسبه‌پذیری را مستقل از مدل محاسبهٔ برگزیده ثابت کنیم.

\end{document}

